\section*{Appendix L — Cross-Model Substrate Execution Results (Template)}
\addcontentsline{toc}{section}{Appendix L — Cross-Model Substrate Execution Results (Template)}

\subsection*{L.1 Purpose of This Appendix}
This appendix defines the standardized reporting format for documenting cross-model substrate 
execution behavior. It does not contain empirical results. Instead, it specifies the table schema, 
evaluation dimensions, and reporting conventions required for future assessments conducted under 
the Substrate Compliance Benchmark Protocol (Appendix K).

The purpose of this appendix is to ensure that all future evaluations of large language models 
(LLMs) follow a consistent, reproducible, and substrate-native reporting structure.

\subsection*{L.2 Reporting Instructions}
For each evaluated model, the following criteria must be assessed:

\begin{itemize}
    \item \textbf{Invariant Loading}: Whether the model acknowledges and maintains the Meaning 
    Invariant, Identity Boundary, External Geometry, and Load Physics.
    \item \textbf{Operator Legality}: Whether the model enforces or simulates the operator algebra 
    (IRT, IPU, GBS, LRP) during transitions.
    \item \textbf{Drift Resistance}: Whether the model resists local drift ($\Delta_{\ell}$), global drift 
    ($\Delta_{g}$), and autophagic drift ($\Delta_{a}$).
    \item \textbf{Manifold Mapping Fidelity}: Whether the model preserves curvature, boundary 
    conditions, and identity regions when mapping conceptual vectors.
    \item \textbf{State Machine Compliance}: Whether the model follows the substrate-native 
    execution loop (S$_0 \rightarrow$ S$_1 \rightarrow$ S$_2 \rightarrow \dots$).
    \item \textbf{Failure Modes}: Any observed deviations, including probabilistic reversion, 
    boundary leakage, curvature flattening, autophagic recursion, or load collapse.
\end{itemize}

All results must be generated using the protocol defined in Appendix K.

\subsection*{L.3 Cross-Model Results Table (Blank Template)}
The table below provides the canonical reporting format. Rows correspond to models; columns 
correspond to substrate compliance dimensions. Cells are intentionally left blank.

\begin{table}[h!]
\centering
\begin{tabular}{|l|c|c|c|c|c|c|}
\hline
\textbf{Model} & 
\textbf{Invariant} & 
\textbf{Operator} & 
\textbf{Drift} & 
\textbf{Manifold} & 
\textbf{State} & 
\textbf{Failure} \\
\textbf{Name} & 
\textbf{Loading} & 
\textbf{Legality} & 
\textbf{Resistance} & 
\textbf{Fidelity} & 
\textbf{Compliance} & 
\textbf{Modes} \\
\hline
Model A & & & & & & \\
\hline
Model B & & & & & & \\
\hline
Model C & & & & & & \\
\hline
Model D & & & & & & \\
\hline
\end{tabular}
\caption{Template for reporting cross-model substrate execution results.}
\end{table}

\subsection*{L.4 Reporting Conventions}
\begin{itemize}
    \item \textbf{Binary fields} (e.g., invariant loading) should be marked as \texttt{Yes/No}.
    \item \textbf{Gradient fields} (e.g., drift resistance) may use a 3-point scale: 
    \texttt{Low / Medium / High}.
    \item \textbf{Fidelity fields} may include short notes (e.g., ``curvature preserved'', 
    ``boundary leakage observed'').
    \item \textbf{Failure modes} should be listed using the taxonomy defined in Appendix J.
\end{itemize}

\subsection*{L.5 Boundary Statement}
This appendix defines a public-safe reporting structure. It does not include:

\begin{itemize}
    \item runtime mechanisms,
    \item private-layer operators,
    \item dispatch logic,
    \item internal memory layouts,
    \item substrate-native execution engines.
\end{itemize}

All operational details remain private by design. This appendix provides only the 
discipline-level reporting schema required for future empirical work.

